\documentclass[aspectratio=169]{beamer}
\usetheme{Madrid}
\usecolortheme{seahorse}

\usepackage[utf8]{inputenc}
\usepackage[T1]{fontenc}
\usepackage{lmodern}
\usepackage{booktabs}
\usepackage{amsmath, amssymb}
\usepackage{array}
\usepackage{hyperref}
\AtBeginSection[]{
    \begin{frame}[plain]
        \centering
        \vfill
        {\Huge\bfseries\insertsectionhead}
        \vfill
    \end{frame}
}

\AtBeginSubsection[]{
    \begin{frame}[plain]
        \centering
        \vfill
        {\Huge\bfseries\insertsubsectionhead}
        \vfill
    \end{frame}
}

\title[EC508 Map]{EC508 Econometrics: Integrated Course Map}
\subtitle{From MLR Foundations to IV, Panel, and Time Series}
\author{Class Summary from slides\_andre}
\date{April 2026}

\begin{document}

\begin{frame}
	\titlepage
\end{frame}

\begin{frame}
	\tableofcontents
\end{frame}


\section{Recap}

\begin{frame}{Basic MLR Recap (Wooldridge)}
Model:
\[
y_i = \beta_0 + \beta_1 x_{i1} + \cdots + \beta_k x_{ik} + u_i
\]
Key assumptions:
\begin{itemize}
	\item MLR.1 Linearity in parameters
	\item MLR.2 Random sampling
	\item MLR.3 No perfect collinearity
	\item MLR.4 Zero conditional mean: $E(u_i\mid X_i)=0$
	\item MLR.5 Homoskedasticity (for usual SE formulas)
\end{itemize}
\end{frame}

\begin{frame}{Why does it matter}
\begin{itemize}
	\item MLR.1--MLR.4 hold: unbiased/consistent
	\item MLR.5 also holds: efficient
	\item If normality is added: exact small-sample $t$ and $F$ inference
\end{itemize}
\vspace{0.3cm}
Typical failure map:
\begin{itemize}
	\item Break MLR.4 $\Rightarrow$ bias/inconsistency (endogeneity)
	\item Break MLR.5 $\Rightarrow$ wrong classical SE/tests, inefficiency
\end{itemize}
\end{frame}

\section{1. Heteroskedasticity}

\begin{frame}{Unit 1: Heteroskedasticity --- What Changes?}
\[
\mathrm{Var}(u_i\mid X_i)=\sigma^2 \quad \text{(homoskedasticity)}
\]
Changed setting:
\[
\mathrm{Var}(u_i\mid X_i)=\sigma_i^2
\]
\begin{itemize}
	\item Unbiased/consistent (if zero conditional mean holds)
	\item Lose efficiency (no longer BLUE)
	\item Lose inference: usual (non-robust) SE, $t$, $F$, confidence intervals are unreliable
\end{itemize}
\begin{block}{Inference fix}
Use heteroskedasticity-robust (Huber-White) standard errors for valid large-sample inference.
\end{block}

\end{frame}

\begin{frame}{Robust Standard Errors}
\begin{itemize}
	\item Consistent under unknown heteroskedasticity forms
	\item $\hat \beta$ unchanged; only estimated uncertainty changes
	\item Robust SE can be larger or smaller; often larger in practice
\end{itemize}
\end{frame}

\begin{frame}{Testing}
\begin{itemize}
	\item Breusch-Pagan (linear variance function idea)
	\item White test (more general nonlinear form using squares/interactions)
\end{itemize}
Logic:
\begin{itemize}
	\item Regress $\hat u_i^2$ on functions of regressors
	\item Significant explanatory power in this auxiliary regression suggests heteroskedasticity
\end{itemize}
\end{frame}

\begin{frame}{WLS/GLS}
If variance form is known, weight:
\[
\mathrm{Var}(u_i\mid X_i)=\sigma^2 h_i
\]
Use weighted least squares:
\begin{itemize}
	\item Lower weight to high-variance observations
	\item Efficiency relative to OLS
	\item Bad variance specification can hurt
\end{itemize}
\end{frame}

\section{2. Data and Specification Issues}

\begin{frame}{Other issues}
\begin{itemize}
	\item Functional form misspecification
	\item Proxy variables
	\item Measurement error (in $y$ vs in $x$)
	\item Missing data and outliers
\end{itemize}
\end{frame}

\begin{frame}{Functional Form}
\begin{itemize}
	\item RESET test
	\item Non-nested comparisons
\end{itemize}
\end{frame}

\begin{frame}{Proxy Variables}
\begin{itemize}
	\item If error is simply random, you are good
	\item If not perfect, usually at least better
	\item In exam, decompose the variable into the proxy and the error
\end{itemize}
\end{frame}


\begin{frame}{Missing Data and Selection}
Cases:
\begin{itemize}
	\item Missing at random: no bias, less precision
	\item Missing based on $Y$ or unobservables: sample selection bias
	\item (In class: draw how can induce attenuation using scatter deleting high income people.)
\end{itemize}
\end{frame}

\begin{frame}{Outliers and Influence}
\begin{itemize}
	\item OLS \textbf{squares} residuals, extreme points have large effect
	\item Report robustness unless a clear data error justifies exclusion
\end{itemize}
\end{frame}

\begin{frame}{Units 1 and 2 Takeaway}
\begin{itemize}
    \item From heteroskedasticity: what breaks, how we solve
    \item From specification/data issues: present more problems that can threaten consistency besides the already known omitted variable approach
\end{itemize}
\end{frame}

\section{3. Time Series}

\begin{frame}{Time Series vs Cross Section}
\begin{itemize}
	\item Ordering matters
	\item Can have seasonality and trends
	\item Random sampling hardly holds
	\item Lag structure, stationarity, trends, and serial dependence
\end{itemize}
\end{frame}

\begin{frame}{Dynamics, interpretation and long-run effects}
Static model:
\[
y_t = \beta_0 + \beta_1 x_t + u_t
\]
Finite distributed lag (FDL):
\[
y_t = \delta_0 + \delta_1 x_t + \delta_2 x_{t-1} + \cdots + u_t
\]
\end{frame}

\begin{frame}{Assumptions: Time Series}
Class assumptions TS.1--TS.5:
\begin{itemize}
	\item TS.1 Linearity in parameters
	\item TS.2 No perfect collinearity
	\item TS.3 Strict exogeneity: $E(u_t\mid X)=0$ (not $X_t$!) for all $t$
	\item TS.4 Homoskedasticity
	\item TS.5 No serial correlation in errors
\end{itemize}
\end{frame}

\begin{frame}{What's new}
\begin{itemize}
    \item Endogeneity is a stronger condition
	\item OLS may remain unbiased under TS.1--TS.3
	\item But efficiency and classical inference rely on TS.4--TS.5
\end{itemize}
\end{frame}

\section{4. Time Series and Forecasting}

\begin{frame}{Autocorrelation and AR Models}
AR(1):
\[
y_t = \beta_0 + \beta_1 y_{t-1} + u_t
\]
AR($p$): $p$ lags.
\begin{itemize}
	\item Prediction vs. causality
	\item Lag choice on information criteria (AIC/BIC)
\end{itemize}
\end{frame}

\begin{frame}{ADL Models and Granger Causality}
ADL combines lags of $y$ and predictors $x$:
\[
y_t = \beta_0 + \sum_{j=1}^p \beta_j y_{t-j} + \sum_{m=1}^r \delta_m x_{t-m} + u_t
\]
Granger causality test: (regretted calling it causality)
\begin{itemize}
	\item Lagged $x$ predicts $y$?
	\item Predictive, not causal
\end{itemize}
\end{frame}

\begin{frame}{Forecasting and Uncertainty}
\begin{itemize}
	\item Forecast error (out-of-sample vs in-sample)
	\item RMSE/RMSFE quantifies typical forecast mistake
	\item Build forecast intervals using variance of errors
	\item In-sample approximation
	\item Pseudo out-of-sample rolling/recursive evaluation
\end{itemize}
\end{frame}

\begin{frame}{Non-Stationarity and Unit Roots}
Random walk with drift:
\[
y_t = \beta_0 + y_{t-1} + u_t
\]
Implication:
\begin{itemize}
	\item Stochastic trend, nonstationary
	\item Regressions in levels can be spurious
	\item Dickey-Fuller test for unit root
\end{itemize}
\end{frame}

\begin{frame}{Unit Roots}
\begin{itemize}
	\item Standard asymptotics fail
	\item DF/ADF test: $$\Delta y_t = \alpha_0 + \alpha_1 y_{t-1} + \sum_{j=1}^p \gamma_j \Delta y_{t-j} + v_t$$
    \item If $\alpha_1=0$: unit root, nonstationary. Need alternative critical values
	\item Differences are stationary (if $I(1)$): $\Delta y_t$
\end{itemize}
\end{frame}

\begin{frame}{Takeaway}
\begin{block}{Where this fits}
    New data structure. Problems with sampling and spurious regressions. 
\end{block}
\end{frame}

\section{5. Panel Data}

\begin{frame}{Data}
    \begin{itemize}
        \item Repeated cross-section and panel data
        \item DGP: \[
y_{it}=\beta_0+\beta_1 x_{it}+v_{it}, \; v_{it}=a_i+u_{it}
\]
        \item POLS: $\mathrm{Cov}(x_{it},a_i)=0$ and $\mathrm{Cov}(x_{it},u_{it})=0$ for all $t$ (if $a_i$ correlated with $x_{it}$, OVB)
    \end{itemize}
\end{frame}

\begin{frame}{First Difference (FD) Estimator}
Two-period illustration:
\[
\Delta y_i = \delta_0 + \beta_1 \Delta x_i + \Delta u_i
\]
What changes:
\begin{itemize}
	\item Differencing removes $a_i$ exactly
	\item $x_i$ must vary
	\item $\Delta u_i$ uncorrelated with $\Delta x_i$
	\item For usual inference: homoskedastic/no problematic serial dependence in differenced errors
\end{itemize}
\end{frame}

\begin{frame}{Fixed Effects (FE / Within) Estimator}
Within transformation:
\[
y_{it}-\bar y_i = \beta_1(x_{it}-\bar x_i) + (u_{it}-\bar u_i)
\]
\begin{itemize}
	\item Removes all time-invariant unobservables $a_i$
	\item Strict exogeneity: $\mathrm{Cov}(x_{it},u_{is})=0$ for all $t,s$
	\item No need for $\mathrm{Cov}(a_i,x_{it})=0$
	\item Equivalent to FD if only two periods but can be more efficient with $T>2$
\end{itemize}
\end{frame}

\begin{frame}{Random Effects (RE) Estimator}
RE keeps $a_i$ in composite error but assumes orthogonality:
\[
\mathrm{Cov}(a_i,x_{it})=0
\]
\begin{itemize}
	\item Bias variance tradeoff: pay the price of this assumption, get higher efficiency
\end{itemize}
\end{frame}

\begin{frame}{Panel Assumptions}
\small
\begin{center}
\begin{tabular}{llll}
\toprule
Estimator & Assumption & Solves $a_i$--$x$ corr. & Time-invariant $x$ \\
\midrule
POLS & $\mathrm{Cov}(x_{it},a_i)=0$ & No & Yes \\
FD & $\mathrm{Cov}(\Delta x_{it},\Delta u_{it})=0$ & Yes & No \\
FE & $\mathrm{Cov}(x_{it},u_{is})=0\;\forall t,s$ & Yes & No \\
RE & $\mathrm{Cov}(x_{it},a_i)=0$ & No & Yes \\
\bottomrule
\end{tabular}
\end{center}
\normalsize
\end{frame}

\begin{frame}{DiD Logic}
Two-period treatment effect representation:
\[
\hat\beta_{DiD} = \Delta \bar y_{treat} - \Delta \bar y_{control}
\]
\end{frame}

\begin{frame}{}
\begin{block}{Where this fits}
Panel methods solve unobserved heterogeneity. Sometimes sacrifice non-time-varying regressors.
\begin{itemize}
	\item If $a_i$ likely correlated with regressors: FE/FD
	\item If orthogonality plausible and time-invariant regressors matter: RE
\end{itemize}
\end{block}
\end{frame}

\section{6. Instrumental Variables}

\begin{frame}{Unit 6: The Endogeneity Problem}
Structural equation:
\[
y = \beta_0 + \beta_1 x + u, \quad \mathrm{Cov}(x,u)\neq 0
\]
\begin{itemize}
	\item Endogeneity (e.g., OVB), biased and inconsistent
	\item Instrument $z$:
\begin{itemize}
	\item Exogeneity (not testable): $\mathrm{Cov}(z,u)=0$
	\item Relevance (testable): $\mathrm{Cov}(z,x)\neq 0$
\end{itemize}
\end{itemize}
$$ \hat \beta = (Z'X)^{-1}Z'y $$ consistent if $z$ is valid
\end{frame}


\begin{frame}{2SLS Workflow}
1st Stage:
\[
x = \pi_0 + \pi_1 z + \pi_2 w + v
\]
2nd Stage:
\[
y = \beta_0 + \beta_1 \hat x + \beta_2 w + e
\]
Key note:
\begin{itemize}
	\item Correct 2SLS standard errors are required
	\item Equivalent to previous one if one instrument, one endogenous
\end{itemize}
\end{frame}

\begin{frame}{Changes}
\begin{itemize}
	\item Consistent, not efficient
	\item Weak instruments can create severe finite-sample problems
\end{itemize}
\end{frame}

\begin{frame}{IV with Multiple Instruments}
Benefits:
\begin{itemize}
	\item More first-stage predictive power can improve precision
	\item Enables overidentification checks when more instruments than endogenous regressors
	\item Endogeneity tests (OLS vs IV logic)
	\item Overidentifying restriction tests (joint instrument exogeneity assumptions)
\end{itemize}
\end{frame}


\begin{frame}{Takeaway}
\begin{block}{Where this fits}
If endogeneity, IV can solve by ``taking exogeneity'' from another variable.
\end{block}
\end{frame}

\section{7. Simultaneous Equations}

\begin{frame}{Simultaneity as Endogeneity}
Example system:
\[
y_1 = \alpha_1 y_2 + \beta_1 z_1 + u_1,\qquad
y_2 = \alpha_2 y_1 + \beta_2 z_2 + u_2
\]
What changes:
\begin{itemize}
	\item Endogenous are jointly determined: regressor in each equation contains part of structural error from the other equation
    \item OLS inconsistent
    \item Need excluded exogenous shifters to identify each structural equation
	\item Order condition (necessary): more excluded exogenous variables than those endogenous (both equations above are identified)
	\item Rank condition (necessary and sufficient): excluded variables shift endogenous regressor in that equation
\end{itemize}
\end{frame}


\begin{frame}{Connection to IV}
\begin{itemize}
	\item Each structural equation with endogenous RHS variable can be estimated by 2SLS if identified
	\item Instruments come from exogenous variables excluded from that equation
\end{itemize}
\end{frame}

\begin{frame}{Unit 7 Takeaway}
\begin{block}{Where this fits}
SEM formalizes equilibrium/reverse-causality environments and shows why identification design is central before estimation.
\end{block}
\end{frame}

\section{Review}

\begin{frame}{Covered}
\begin{enumerate}
	\item Stop assuming homoskedasticity
	\item Explore additional ways to break MLR.4 (endogeneity)
	\item New data type: time series
	\item New data type: panel data
	\item A solution to endogeneity: instrumental variables
	\item Simultaneous equations
\end{enumerate}
\end{frame}


\end{document}
